\section*{Chapitre \ref{ch:b2:phanerozoic-earth} --- La Terre phanérozoïque}

\begin{solution}{exo:b2:phanerozoic-earth:1}
(a) Paris (\ang{5.1}~N) et New York (\ang{3.6}~S)~; Le Cap et Sydney se trouvaient près de \ang{60}~S.
(b) Le Cap~: $61.9 - 33.9 = \ang{28.0}$, $28.0 \times 111.2 = \qty{3114}{km}$~; New York~:
$40.7 + 3.6 = \ang{44.3}$, \qty{4926}{km}. (c) $3114/300 = \qty{10.4}{km}$ par million d'années,
soit \qty{1.0}{cm} par an~; $4926/300 = \qty{16.4}{km}$ par million d'années, soit \qty{1.6}{cm} par an.
\end{solution}

\begin{solution}{exo:b2:phanerozoic-earth:2}
(a) 72, 120.1, 50.5 et \qty{135.4}{Ma}~; moyenne $(444 - 66)/4 = \qty{94.5}{Ma}$. (b) La fin du
Permien (Sibérie, \qty{1.9}{Ma} d'écart), la fin du Trias (CAMP, \qty{0.4}{Ma}) et la fin du Crétacé
(Deccan, \qty{0.5}{Ma}). (c) Non~: les intervalles vont de 50 à \qty{135}{Ma}~; une moyenne existe
pour tout ensemble de dates.
\end{solution}

\begin{solution}{exo:b2:phanerozoic-earth:3}
(a) $\qty{91}{m}/\qty{0.036}{mm.an^{-1}} = \num{2.5e6}$ années. (b) $3.7/0.036 = 103$ fois. (c)
Loin de tout rivage, en eau claire, peu de vase venue des fleuves atteignait le fond~; la pluie
de coquilles d'algues était le principal sédiment.
\end{solution}

\begin{solution}{exo:b2:phanerozoic-earth:4}
(a) Du grès fluviatile jusqu'à la craie, les dépôts s'approfondissent vers le haut~: c'est une
transgression. Le sable de plage au-dessus de la craie enregistre une régression. (b) Non~: la coupe
enregistre une hausse du niveau relatif de la mer, que peuvent aussi produire l'affaissement du lieu ou une baisse de l'apport de sédiments. (c)
Trouver la même transgression, au même moment, sur d'autres continents, ou retirer l'affaissement
par \omterm{def:b2:basins-compared:backstripping}{backstripping}.
\end{solution}

\begin{solution}{exo:b2:phanerozoic-earth:5}
(a) $233/97 = \qty{2.4}{m}$ par million d'années, soit \qty{2.4}{\micro m} par an. (b) $3.7/0.0024 =
1540$~: la hausse actuelle est environ 1500 fois plus rapide. (c) La courbe corrigée de l'affaissement~: les sondages
ont montré que les courbes des cycles donnent bien la date des baisses mais surestiment leur amplitude d'un
facteur deux au moins~; New Jersey a son propre biais (topographie dynamique), d'où l'estimation de $100 \pm
\qty{50}{m}$.
\end{solution}

\begin{solution}{exo:b2:phanerozoic-earth:6}
(a) $\num{e5} \times \num{e9} \times 2900 = \qty{2.9e17}{kg}$ de basalte, et $\num{2.9e17}
\times 0.0075 = \qty{2.2e15}{kg}$ de dioxyde de carbone, $\num{2.2e15}/0.04401 = \qty{4.9e16}{mol}$.
(b) $\num{278e-6} \times (44.01/28.97) \times \num{5.1e18} = \qty{2.15e15}{kg}$. (c) Une pulsation
a libéré à peu près autant que l'ensemble de l'air de 1750. (d) $\num{e17}/\num{2.15e15} = 46.5$~: près de cinquante.
\end{solution}

\begin{solution}{exo:b2:phanerozoic-earth:7}
(a) $2630 \times \frac{\pi}{6} \times (\num{e4})^3 = \qty{1.38e15}{kg}$. (b) $\tfrac12 \times
\num{1.38e15} \times (\num{11.2e3})^2 = \qty{8.7e22}{J}$~; à \qty{20}{km.s^{-1}},
\qty{2.76e23}{J}. (c) Une année de \qty{47}{TW} vaut $\num{47e12} \times \num{3.156e7} =
\qty{1.48e21}{J}$~: 59 et 186 ans.
\end{solution}

\begin{solution}{exo:b2:phanerozoic-earth:8}
(a) Couleur sombre, lamination fine, absence de terriers, beaucoup de matière organique et de pyrite~; souvent
un décalage des isotopes du carbone et des signes de réchauffement. (b) L'événement du Toarcien avec le Karoo--Ferrar, l'EAO~1a avec
Ontong Java. (c) Les dates portent des incertitudes pouvant atteindre des millions d'années, et il existe beaucoup
d'événements et de provinces, de sorte que des coïncidences fortuites sont probables~; un lien causal exige que la province commence
juste avant l'événement et un traceur du volcanisme, tel le mercure, dans les mêmes sédiments.
\end{solution}

\begin{solution}{exo:b2:phanerozoic-earth:9}
(a) $60.9 + 5.1 = \ang{66.0}$, $66.0 \times 111.2 = \qty{7339}{km}$ en \qty{200}{Ma}~:
\qty{36.7}{km} par million d'années, \qty{3.7}{cm} par an. (b) Plus du double des
\qty{1.6}{cm} par an depuis \qty{300}{Ma}. (c) Seul le changement de latitude est connu~; le changement de
longitude, libre dans un référentiel paléomagnétique, et tout détour s'ajoutent à la distance.
\end{solution}

\begin{solution}{exo:b2:phanerozoic-earth:10}
(a) La fraction du plancher d'âge $t$ vaut $p(t) = \frac2T(1 - t/T)$. Alors $\int_0^T tp\,\mathrm{d}t =
\frac2T\bigl(\frac{T^2}2 - \frac{T^2}3\bigr) = \frac T3$ et $\int_0^T \sqrt t\,p\,\mathrm{d}t =
\frac2T\bigl(\frac23 T^{3/2} - \frac25 T^{3/2}\bigr) = \frac8{15}\sqrt T$. (b) $T = 3 \times 64.2 =
\qty{192.6}{Ma}$~; $2500 + \frac8{15} \times 350 \times 13.878 = \qty{5091}{m}$. (c) $T/2 =
\qty{96.3}{Ma}$~: \qty{4332}{m}~; la baisse est de \qty{759}{m}, \qty{2}{\percent} de moins que
\qty{774}{m}. (d) Modèle réparti~: $\frac23\sqrt\tau = \frac23\sqrt{2\bar t} = 0.943\sqrt{\bar t}$~;
triangle~: $\frac8{15}\sqrt{3\bar t} = 0.924\sqrt{\bar t}$. La réponse dépend bien plus de l'âge moyen
que de la forme.
\end{solution}

\begin{solution}{exo:b2:phanerozoic-earth:11}
(a) Dans $\Delta s$, seul le facteur $1 - f^{-1/2}$ dépend de $f$~; diviser par sa valeur en $f = 2$,
$1 - 2^{-1/2} = 0.2929$. (b) $f = (1 - 0.2929\,\Delta s/432)^{-2}$~: 1.071 pour \qty{50}{m}, 1.239
pour \qty{150}{m}~; soit de 7 à \qty{24}{\percent} plus vite qu'aujourd'hui. (c) $432 \times (1 -
0.8^{-1/2})/0.2929 = \qty{-174}{m}$~: la mer baisserait de \qty{174}{m}. (d) La répartition des âges
ne change qu'au fur et à mesure que le vieux plancher est subduit et remplacé, ce qui prend un temps de l'ordre de $\tau$,
soit environ \qty{100}{Ma}~; la mer suit sur des dizaines de millions d'années.
\end{solution}

\begin{solution}{exo:b2:phanerozoic-earth:12}
(a) $2^{-7.4/2} = 0.077$~: \qty{92}{\percent} disparus. (b) $t = 2\log_2(1/0.19) = \qty{4.8}{Ma}$.
(c) $4.8/0.060 = 80$ fois plus vite. (d) À l'intérieur d'un long étage, la seule extinction de fond fait disparaître
la plupart des espèces, de sorte qu'une courte pulsation s'ajoute à un total déjà élevé et se distingue mal de lui~; la
lacune de l'enregistrement décale aussi vers le passé les dernières apparitions.
\end{solution}

\begin{solution}{pb:b2:phanerozoic-earth:1}
\textbf{1.} 2500, 4250, 5304 et \qty{6000}{m}.
\textbf{2.} $\tau = 2 \times 64.2 = \qty{128.4}{Ma}$~; $2500 + 350\sqrt{128.4} = \qty{6466}{m}$.
\textbf{3.} $\num{293e6}/(65000 \times 128.4) = \qty{35.1}{km}$ par million d'années,
\qty{35.1}{mm} par an~: une vitesse lente.
\textbf{4.} $\frac1\tau\int_0^\tau(2500 + 350\sqrt t)\,\mathrm{d}t = 2500 + \frac23 \times
350\sqrt\tau = \qty{5144}{m}$.
\textbf{5.} La loi en racine carrée est trop profonde pour le plancher de plus de quelques dizaines de millions d'années,
qui s'aplatit~; les âges réels ne sont pas répartis également, le plancher jeune étant le plus abondant~; et
les sédiments, \omterm{def:b2:oceanic-islands:seamount}{monts sous-marins} et plateaux relèvent le plancher.
\textbf{6.} $\num{293e6} \times 5.144 = \qty{1.51e9}{km^3}$, \qty{13}{\percent} de plus que
l'océan entier~: le modèle est trop profond.
\textbf{7.} $v = \qty{17.6}{km}$ par million d'années~; $17.6 \times 128.4 = \qty{2260}{km}$ de
chaque côté.
\textbf{8.} $t = 1000/17.6 = \qty{56.8}{Ma}$~; $2500 + 350\sqrt{56.8} = \qty{5138}{m}$.
\textbf{9.} D'un seul flanc, avec $x = vt$~: $\int_0^{v\tau}\bigl(d(\tau) - d(x/v)\bigr)\mathrm{d}x
= vc\int_0^\tau(\sqrt\tau - \sqrt t)\,\mathrm{d}t = \frac13 vc\,\tau^{3/2}$~; le double pour les deux.
$\frac23 \times 17.6 \times 0.350 \times 128.4^{3/2} = \frac23 \times 17.6 \times 0.350 \times
1454.9 = \qty{5975}{km^2}$.
\textbf{10.} $65000 \times 5975 = \qty{3.88e8}{km^3}$.
\textbf{11.} $\num{3.88e8}/\num{1.33e9} = 0.29$.
\textbf{12.} Son plancher est plus jeune en moyenne, donc plus chaud, moins dense et plus haut~; la bande de
plancher jeune le long de l'axe est plus large.
\textbf{13.} $\tau = \qty{64.2}{Ma}$~; $2500 + \frac23 \times 350\sqrt{64.2} = \qty{4370}{m}$.
\textbf{14.} $5144 - 4370 = \qty{774}{m}$.
\textbf{15.} $\num{293e6} \times 0.774 = \qty{2.27e8}{km^3}$.
\textbf{16.} $\num{2.27e8}/\num{362e6} = \qty{0.627}{km}$, soit \qty{627}{m}.
\textbf{17.} $\Delta s = 627 \times (3300 - 1025)/3300 = \qty{432}{m}$.
\textbf{18.} $f = \bigl(1 - 0.2929 \times 100/432\bigr)^{-2} = 1.15$~: \qty{15}{\percent} plus vite.
\textbf{19.} $0.149 \times 148 = \qty{22.1}{\mega km^2}$.
\textbf{20.} $48.1 + 0.32 \times 7.7 = \qty{50.6}{\percent}$, soit \qty{74.9}{\mega km^2}.
\textbf{21.} En montant, la mer couvre une aire croissante, de sorte que le même volume forme une couche plus mince
(\cref{prop:b2:phanerozoic-earth:flooding}).
\textbf{22.} Aires submergées à 100, 200, 300 et \qty{400}{m}~: 22.1, 41.7, 57.3 et
\qty{71.2}{\mega km^2}. Paliers~: $(362 + 11.0) \times 0.1 = \qty{37.3e6}{km^3}$, puis 39.4, 41.2
et \qty{42.6e6}{km^3}~; cumulés 37.3, 76.7, 117.9 et \qty{160.5e6}{km^3}. Le volume à placer
est \qty{156.4e6}{km^3}, entre 300 et \qty{400}{m}~: $300 + (156.4 - 117.9)/42.6 \times 100
= \qty{390}{m}$.
\textbf{23.} $38.7 + 0.90 \times 9.4 = \qty{47.2}{\percent}$ des terres actuelles.
\textbf{24.} Doubler la vitesse d'expansion élèverait la mer d'environ \textbf{\qty{390}{m}} et submergerait près de la moitié
des terres actuelles.
\end{solution}
